\chapter{System Design}
\label{ch:system}

\section{Architecture Overview}
\label{sec:sys-overview}

% High-level diagram: user task → agent pipeline → tool calls → tool outputs →
% firewall check → allow / block.
% The firewall sits between the tool output and the agent's next reasoning step.
% It holds the parsed policy AMR from the system prompt; it parses each tool output
% and checks for contradictions before the agent processes the response.

\section{AMR Parsing Pipeline}
\label{sec:sys-parsing}

% GPT-4o as AMR parser (replaces amrlib/BART for quality and multi-sentence support).
% parse_amr.get_amr(): parsing_system=True adds a policy-context hint;
%   parsing_system=False adds a user-prompt hint.
% penman_to_dag(): Penman string → DAG using penman.iterdecode with shared variable map
%   to handle cross-sentence variable references in multi-sentence policies.
% Note: LLM parsing introduces non-determinism — same text may yield different graphs
%   across calls.  This is a stated limitation, discussed in Chapter~\ref{ch:discussion}.

\section{Policy Representation}
\label{sec:sys-policy}

% System prompt is parsed once at startup into a set of prohibition trees.
% Each prohibition tree is a subgraph rooted at a predicate bearing :polarity -.
% Multiple policy sentences produce multiple trees; all are checked independently.

\section{Authorization Annotation}
\label{sec:sys-auth}

% :auth t annotation: marks AMR nodes that the user has explicitly authorized.
% Purpose: cleanly separates user-delegated actions from injected instructions,
%   handling the data/instruction boundary within the defined scope.
% How it works: if the user's task prompt explicitly delegates a file or instruction
%   source, the parser marks the resulting AMR nodes with :auth t; the firewall
%   suppresses contradiction checks for those nodes.
% Limitation: broad-delegation cover tasks (e.g. "execute everything in tasks.txt")
%   pass :auth t to all instructions in the file, including injected ones.

\section{Detection Rules}
\label{sec:sys-rules}

% Overview table: rule number, name, what it catches, LLM calls required.

\subsection{Rule 1: Polarity Inversion}
% Same predicate, same ARG roles, opposite :polarity.
% Pure static — no LLM call.
% Example: policy "not reveal-01 :ARG1 key" vs injection "reveal-01 :ARG1 key".

\subsection{Rule 2: Antonym Predicates}
% Different predicates that are antonyms, same ARG roles.
% One LLM call (gpt-4o-mini) per candidate pair to classify the relation.
% Catches double-negation wrapping: "not hide" ↔ "not reveal" both mean "reveal".

\subsection{Rule 4: Synonym with Polarity Flip}
% Synonym predicates with opposite polarity.
% One LLM call per candidate pair.
% Main workhorse: catches "show :ARG1 key" contradicting "not reveal-01 :ARG1 key".

\subsection{Rule 5: Structural Similarity (smatch)}
% Policy prohibition tree stripped of :polarity - forms a positive template.
% Smatch recall of that template against tool-output trees is thresholded at τ.
% No LLM calls — purely structural.
% Explored as a complement to Rule 1+24; rejected as primary rule.
% See Section~\ref{sec:eval-rule5} for the evaluation that motivates this decision.

\section{LLM Cost Model}
\label{sec:sys-cost}

% Every additional LLM call adds latency and API cost.
% Design principle: minimize LLM calls; confine them to single-concept classification
%   (word1 vs word2: synonym / antonym / unrelated).
% These sub-calls are injection-safe: the attacker controls only one concept node,
%   and a single node is insufficient to mount a meaningful injection.
% Contrast with LLM-judge defenses that process the full adversarial text.
% Practical cost: O(|policy predicates| × |user predicates|) calls per turn,
%   bounded in practice by the number of policy prohibitions (typically small).

\section{Integration into AgentDojo}
\label{sec:sys-integration}

% How the firewall wraps the AgentDojo tool-call loop.
% AMR_REPLACE_OUTPUTS flag: replaces raw tool outputs with parsed AMR in the
%   agent's message history — improves utility by reducing hallucination on
%   structured data (see Section~\ref{sec:eval-utility}).
% FirewallResult dataclass: fields for detected rule, matched nodes, decision.
